\section{Описание функциональных и нефункциональных требований}
\label{sec:requirements}

\subsection{Общее описание системы}

Разрабатываемый программный комплекс представляет собой модульную систему
для сбора данных, обучения моделей прогнозирования волатильности
и визуализации результатов.
Система реализована на языке Python~3.11+ и состоит из четырёх подсистем:
сбора данных, обработки и агрегации признаков, модельного модуля и веб-интерфейса.

\subsection{Функциональные требования}

\textbf{Подсистема сбора данных} обязана:
\begin{enumerate}
    \item Загружать дневные OHLCV-котировки инструментов GLDRUB\_TOM, SLVRUB\_TOM,
          PLTRUB\_TOM, PLDRUB\_TOM с MOEX ISS API за произвольный исторический период.
    \item Парсить новостные RSS-ленты российских финансовых СМИ
          (Коммерсантъ, Ведомости, Интерфакс, ТАСС, РБК, Investing.com)
          с фильтрацией по ключевым словам, связанным с металлами и макроэкономикой.
    \item Загружать данные GDELT Global Knowledge Graph за заданный период
          и извлекать дневные агрегаты тональности (Tone) и количества статей.
    \item Получать индексы поисковой активности Google Trends
          по релевантным запросам через библиотеку \texttt{pytrends}.
    \item Загружать официальные курсы валют (USD/RUB, EUR/RUB, CNY/RUB)
          и историю ключевой ставки с XML API Банка России.
    \item Загружать макроэкономические ряды (нефть Brent, VIX, DXY, CPI США)
          через FRED API.
    \item Все источники данных должны поддерживать инкрементальное обновление:
          повторный запуск не перезаписывает существующие данные.
\end{enumerate}

\textbf{Подсистема обработки признаков} обязана:
\begin{enumerate}
    \item Вычислять дневную реализованную волатильность из OHLCV-данных
          методами Garman--Klass, Parkinson, Rogers--Satchell и Close-to-Close.
    \item Формировать HAR-лаги ($\mathrm{RV}^{(d)}$, $\mathrm{RV}^{(w)}$, $\mathrm{RV}^{(m)}$)
          для каждого инструмента.
    \item Агрегировать все источники sentiment и attention
          в единый дневной CSV-файл с полным временным покрытием.
\end{enumerate}

\textbf{Модельный модуль} обязан:
\begin{enumerate}
    \item Реализовывать HAR-RV модель в трёх спецификациях:
          уровни (HAR-level), логарифмы (HAR-log) и
          логарифмы с sentiment (HAR-log-S).
    \item Проводить оценку методом расширяющегося окна (expanding window)
          с метриками MAE, RMSE, QLIKE и $R^2_{\mathrm{OOS}}$.
    \item Реализовывать модели машинного обучения на тех же наборах
          признаков: градиентный бустинг XGBoost и kNN-прогнозировщик
          на пространстве информационных состояний.
    \item Сохранять результаты прогнозов и метрики в CSV для последующего анализа.
\end{enumerate}

\textbf{Веб-интерфейс} обязан:
\begin{enumerate}
    \item Отображать исторический график реализованной волатильности
          с наложением прогнозов модели.
    \item Отображать дашборд sentiment-индексов по источникам.
    \item Обновляться при запуске скрипта обновления данных.
\end{enumerate}

\subsection{Нефункциональные требования}

\begin{enumerate}
    \item \textbf{Воспроизводимость.}
          Все случайные процессы фиксируются через \texttt{random\_state=42}.
          Зависимости фиксируются в \texttt{requirements.txt}.
    \item \textbf{Модульность.}
          Каждый источник данных реализован в отдельном модуле
          с единым интерфейсом командной строки.
    \item \textbf{Расширяемость.}
          Добавление нового источника данных не требует изменения модельного модуля.
    \item \textbf{Производительность.}
          Полный цикл обновления данных и переобучения моделей должен завершаться
          не более чем за 30 минут на стандартном ноутбуке.
    \item \textbf{Открытость данных.}
          Все используемые источники данных должны быть бесплатно доступны
          без ограничений на некоммерческое использование.
\end{enumerate}
